\section{Elliptical Priors and Small-Budget Index Decompositions}
\label{sec:elliptical-decompositions-appendix}

The previous section developed the optimal small-budget design in terms of small-budget indices under an arbitrary prior.
In this section we show that the small-budget indices can be written in a more interpretable form under an elliptical prior.
Only this section restricts the prior to the class of elliptical distributions; the general small-budget and large-budget results do not.

We refer to \citeauthor{FangKo90}'s \citetitle{FangKo90}~\cite{FangKo90} for the definition of elliptical distributions and an extensive discussion of their properties.
Note that when $\mathrm{EC}_D(\bm\mu,\Sigma,\phi)$ has finite second moments, the location $\bm\mu$ and scale $\Sigma$ parameters have the following relationship to the mean and covariance (Theorem 2.17 of~\cite{FangKo90}):
\[
  \E[\bm\tau]=\bm\mu \quad\text{and}\quad \cov(\bm\tau) \propto \Sigma.
\]

The key property of elliptical distributions $\mathrm{EC}_D(\bm\mu,\Sigma,\phi)$ that we will use is that $\bm\tau$ has an affine conditional mean given any linear summary of $\bm\tau$.

\begin{lemma}[Linear-summary conditional mean]
  \label{lem:elliptical-conditional-mean-singular-appendix}
  Let \(\bm\tau\sim\mathrm{EC}_D(\bm\mu,\Sigma,\phi)\) have finite second moments.
  Let \(A\in\mathbb R^{D\times k}\), and set \(L=A^\top\bm\tau\).
  Then
  \[
    \E[\bm\tau\mid L]
    =
    \bm\mu
    +
    \cov(\bm\tau,L)\var(L)^\dagger(L-\E[L])
    \qquad\text{a.s.},
  \]
  where \({}^\dagger\) denotes the Moore--Penrose inverse.
\end{lemma}

\begin{proof}
  Let
  \[
    \bm X\coloneqq \bm\tau-\bm\mu,
    \qquad
    C\coloneqq \cov(\bm\tau).
  \]
  By Theorem 2.17 of~\cite{FangKo90}, \(\E[\bm\tau]=\bm\mu\) and \(C\propto\Sigma\).
  If \(C=0\), then
  \(\bm\tau=\bm\mu\) a.s., and the claim is immediate.

  Hence assume \(C\neq 0\).
  Let \(q=\operatorname{rank}(\Sigma)\).
  Then \(q>0\), and \(C=c\Sigma\) for some \(c>0\).
  By Corollary 1 to Theorem 2.14 of~\cite{FangKo90}, there exist
  \(\Gamma\in\mathbb R^{q\times D}\), a nonnegative random variable \(r\),
  and \(\bm u^{(q)}\) uniformly distributed on the unit sphere in
  \(\mathbb R^q\), with \(r\) independent of \(\bm u^{(q)}\), such that
  \[
    \Gamma^\top\Gamma=\Sigma,
    \qquad
    \Gamma^\top(r\bm u^{(q)})\stackrel d=\bm X .
  \]
  Since the desired conditional-mean identity is determined by the joint
  law of \((\bm X,A^\top\bm X)\), it suffices to prove it for this equal-in-distribution copy.
  Thus, writing
  \[
    \bm S\coloneqq r\bm u^{(q)},
  \]
  we may assume, without loss of generality, that
  \[
    \bm X=\Gamma^\top\bm S.
  \]
  The finite-second-moment assumption ensures \(\bm S\) is integrable.

  We first record a simple spherical projection identity.
  For any \(M\in\mathbb R^{q\times k}\), let
  \[
    P\coloneqq M(M^\top M)^\dagger M^\top .
  \]
  Then \(P\) is the orthogonal projection onto \(\operatorname{Im}(M)\), and
  \[
    \E[\bm S\mid M^\top\bm S]=P\bm S
    \qquad\text{a.s.}
  \]
  Indeed, \(P\bm S=M(M^\top M)^\dagger M^\top\bm S\) is measurable with respect to \(M^\top\bm S\).
  Let \(J=2P-I_q\).
  Then \(J\) is orthogonal,
  \(M^\top J\bm S=M^\top\bm S\), and
  \((I_q-P)J\bm S=-(I_q-P)\bm S\).
  Since \(\bm u^{(q)}\) is uniform on the sphere, \(J\bm S\stackrel d=\bm S\).
  Thus, for every bounded measurable
  \(g\),
  \[
    \begin{aligned}
      \E\!\left[g(M^\top\bm S)(I_q-P)\bm S\right]
      &=
      \E\!\left[g(M^\top J\bm S)(I_q-P)J\bm S\right] \\
      &=
      -\E\!\left[g(M^\top\bm S)(I_q-P)\bm S\right].
    \end{aligned}
  \]
  Hence \(\E[(I_q-P)\bm S\mid M^\top\bm S]=0\), proving the identity.

  Apply this identity with \(M=\Gamma A\).
  Since
  \(L=A^\top\bm\mu+A^\top\bm X\), conditioning on \(L=A^\top\bm X\) is the same as
  conditioning on \(A^\top\bm X=M^\top\bm S\).
  Therefore
  \[
    \begin{aligned}
      \E[\bm X\mid L]
      &=
      \Gamma^\top\E[\bm S\mid M^\top\bm S] \\
      &=
      \Gamma^\top M(M^\top M)^\dagger M^\top\bm S \\
      &=
      \Sigma A(A^\top\Sigma A)^\dagger A^\top\bm X .
    \end{aligned}
  \]
  Since \(C=c\Sigma\),
  \[
    \Sigma A(A^\top\Sigma A)^\dagger
    =
    C A(A^\top C A)^\dagger .
  \]
  Finally,
  \[
    \cov(\bm\tau,L)=CA,
    \qquad
    \var(L)=A^\top C A,
    \qquad
    L-\E[L]=A^\top\bm X.
  \]
  Substitution gives the claim.
\end{proof}

For the single-stage objectives, the general-prior formulas from \Cref{prop:first-stage-coefficient-formulas-appendix} already give
\[
  \theta_d^{\singlesuccess}=\E[\tau_d]=\mu_d,
  \qquad
  \theta_d^{\singleimpact}
  =
  \E[\ATE_3] \, \mu_d+\cov(\tau_d,\ATE_3).
\]
The next lemma records the elliptical-prior calculation that will be used for the two-stage objectives.

\begin{lemma}[Weighted linear-summary decomposition]
  \label{lem:elliptical-weighted-linear-summary-appendix}
  Let \(\bm\tau\sim\mathrm{EC}_D(\bm\mu,\Sigma,\phi)\) have finite second moments.
  Let \(L=A^\top\bm\tau\in\mathbb R^k\) be any linear summary.
  Let \(H\) be real-valued and measurable with respect to \(L\), and suppose
  \[
    \E[|H|]<\infty
    \qquad\text{and}\qquad
    \E[|H|\,\|\bm\tau\|_2]<\infty .
  \]
  Then, for every type \(d\),
  \[
    \E[\tau_d H]
    =
    \E[H]\,\mu_d
    +
    \cov(\tau_d,L)^\top\var(L)^\dagger\cov(L,H).
  \]
\end{lemma}

\begin{proof}
  By \Cref{lem:elliptical-conditional-mean-singular-appendix},
  \[
    \E[\bm\tau\mid L]
    =
    \bm\mu
    +
    \cov(\bm\tau,L)\var(L)^\dagger(L-\E[L])
    \qquad\text{a.s.}
  \]
  Taking the \(d\)-th coordinate and using that \(H\) is measurable with
  respect to \(L\),
  \[
    \begin{aligned}
      \E[\tau_d H]
      &=
      \E\!\left[H\,\E[\tau_d\mid L]\right] \\
      &=
      \E[H]\,\mu_d
      +
      \cov(\tau_d,L)^\top\var(L)^\dagger
      \E\!\left[(L-\E[L])H\right] \\
      &=
      \E[H]\,\mu_d
      +
      \cov(\tau_d,L)^\top\var(L)^\dagger\cov(L,H).
    \end{aligned}
  \]
\end{proof}

We now use this result to decompose the small-budget indices for the two-stage success objective and the pilot (two-stage) impact objective into more interpretable components.
We begin with the two-stage success objective.

\begin{proposition}[Two-stage success coefficient under an elliptical prior]
  \label{prop:two-stage-success-elliptical-appendix}
  Assume the prior distribution of \(\bm\tau\) is elliptical with finite second moments.
  Define
  \[
    W_2^{\twosuccess}
    \defeq
    \begin{cases}
      \dfrac{\cov(\ATE_2,p_2(\ATE_2))}{\var(\ATE_2)}, & \text{if }\var(\ATE_2)>0,\\[8pt]
      0, & \text{if }\var(\ATE_2)=0.
    \end{cases}
  \]
  Then, for every type \(d\),
  \[
    \theta_d^{\twosuccess}
    =
    \P(\Test_2=1)\,\mu_d
    +
    W_2^{\twosuccess}\,\cov(\tau_d,\ATE_2),
  \]
  and consequently
  \[
    \pi_d^{\twosuccess}
    =
    \frac{
      \P(\Test_2=1)\,\E[\tau_d]
      +
      W_2^{\twosuccess}\,\cov(\tau_d,\ATE_2)
    }{\sqrt{c_d}}.
  \]
  Moreover, if \(\var(\ATE_2)>0\), then \(W_2^{\twosuccess}>0\).
\end{proposition}

\begin{proof}
  Apply \Cref{lem:elliptical-weighted-linear-summary-appendix} with \(L=\ATE_2\) and \(H=p_2(\ATE_2)\).
  Since \(\E[p_2(\ATE_2)]=\P(\Test_2=1)\), this gives the displayed decomposition.
  When \(\var(\ATE_2)=0\), the covariance term is identically zero and the displayed definition sets \(W_2^{\twosuccess}=0\).
  When \(\var(\ATE_2)>0\), the generalized inverse reduces to ordinary division, giving the displayed ratio.

  It remains to prove positivity in the nondegenerate case.
  Let \(X\) and \(X'\) be independent copies of \(\ATE_2\).
  Since
  \[
    p_2(x)=\Phi(\sqrt{n_2}x-z_{1-\alpha_2})
  \]
  is strictly increasing,
  \[
    (X-X')\bigl(p_2(X)-p_2(X')\bigr)\ge0
    \qquad\text{a.s.},
  \]
  and the inequality is strict on \(\{X\ne X'\}\).
  If \(\var(\ATE_2)>0\), then \(\P(X\ne X')>0\).
  Therefore
  \[
    2\cov(\ATE_2,p_2(\ATE_2))
    =
    \E\!\left[(X-X')\bigl(p_2(X)-p_2(X')\bigr)\right]
    >0.
  \]
  Dividing by \(\var(\ATE_2)>0\) proves \(W_2^{\twosuccess}>0\).
\end{proof}

Now we turn to the main result of this section, the decomposition of the small-budget index for the pilot (two-stage) impact objective.

\refstepcounter{proposition}
\begin{linkedresult}{proof}{prop:marginal-impact-interpretation}
  \label{prop:marginal-impact-interpretation-appendix}
  Assume the prior distribution of \(\bm\tau\) is elliptical with finite second moments.
  Then there exist scalars \(W_1,W_2,W_3\in\R\), not depending on \(d\), such that, for every type \(d\),
  \[
    \theta_d
    =
    W_1\,\mu_d
    +
    W_2\,\cov(\tau_d,\ATE_2)
    +
    W_3\,\cov(\tau_d,\ATE_3).
  \]
  Consequently,
  \[
    \idx_d
    =
    \frac{
      W_1\,\E[\tau_d]
      +
      W_2\,\cov(\tau_d,\ATE_2)
      +
      W_3\,\cov(\tau_d,\ATE_3)
    }{\sqrt{c_d}}.
  \]

  If additionally \(\bm s_2\propto\bm s_3\), \(\E[\ATE_3]\ge0\), and \(\var(\ATE_3)>0\), then
  \[
    \theta_d
    \propto
    \overline W_1\,\mu_d+
    \cov(\tau_d,\ATE_3)
  \]
  for some \(\overline W_1>0\).
  Equivalently, the small-budget index is a positive multiple of
  \[
    \idx_d
    \propto
    \frac{
      \overline W_1\,\E[\tau_d]
      +
      \cov(\tau_d,\ATE_3)
    }{\sqrt{c_d}},
    \qquad
    \overline W_1>0.
  \]
\end{linkedresult}

\begin{proof}
  Let
  \[
    L\defeq
    \begin{pmatrix}
      \ATE_2\\
      \ATE_3
    \end{pmatrix},
    \qquad
    H\defeq p_2(\ATE_2)\ATE_3.
  \]
  Since \(0<p_2(\ATE_2)<1\) and \(\ATE_3\) is a linear summary of \(\bm\tau\), finite second moments imply \(\E[|H|\,\|\bm\tau\|_2]<\infty\).
  Applying \Cref{lem:elliptical-weighted-linear-summary-appendix} gives
  \[
    \theta_d
    =
    \E[\tau_dH]
    =
    \E[H] \, \mu_d
    +
    \cov(\tau_d,L)^\top\var(L)^\dagger\cov(L,H).
  \]
  Define
  \[
    W_1\defeq \E[H],
    \qquad
    \begin{pmatrix}
      W_2\\
      W_3
    \end{pmatrix}
    \defeq
    \var(L)^\dagger\cov(L,H).
  \]
  Since
  \[
    \cov(\tau_d,L)
    =
    \begin{pmatrix}
      \cov(\tau_d,\ATE_2)\\
      \cov(\tau_d,\ATE_3)
    \end{pmatrix},
  \]
  this is exactly
  \[
    \theta_d
    =
    W_1\mu_d
    +
    W_2\cov(\tau_d,\ATE_2)
    +
    W_3\cov(\tau_d,\ATE_3).
  \]
  Dividing by \(\sqrt{c_d}\) gives the displayed index decomposition.

  It remains to prove the representative-follow-up reduction.
  Suppose \(\bm s_2\propto\bm s_3\), which implies \(\ATE_2=\ATE_3\) a.s.
  Write their common value as
  \[
    X\defeq \ATE_3.
  \]
  Then
  \[
    \theta_d
    =
    \E[\tau_d p_2(X)X].
  \]
  Since \(\var(X)=\var(\ATE_3)>0\), applying \Cref{lem:elliptical-weighted-linear-summary-appendix} with \(L=X\) and \(H=p_2(X)X\) gives
  \[
    \theta_d
    =
    A_0\mu_d+B_0\cov(\tau_d,X),
  \]
  where
  \[
    A_0\defeq \E[p_2(X)X],
    \qquad
    B_0\defeq \frac{\cov(X,p_2(X)X)}{\var(X)}.
  \]
  We show that \(A_0>0\) and \(B_0>0\) when \(\E[X]\ge0\).

  Let \(m\defeq\E[X]\).
  As in the proof of \Cref{prop:two-stage-success-elliptical-appendix}, strict monotonicity of \(p_2\) and \(\var(X)>0\) imply
  \[
    \cov(X,p_2(X))>0.
  \]
  Hence
  \[
    A_0
    =
    \E[p_2(X)]\,m+
    \cov(X,p_2(X))
    >0
  \]
  whenever \(m\ge0\).
  Also,
  \begin{align*}
    \cov(X,p_2(X)X)
    &=
    \E[(X-m)p_2(X)X] \\
    &=
    m\,\cov(X,p_2(X))
    +
    \E[p_2(X)(X-m)^2]
    >0,
  \end{align*}
  because \(m\ge0\), \(\cov(X,p_2(X))>0\), \(p_2(X)>0\) a.s., and \(\var(X)>0\).
  Thus \(B_0>0\).
  Therefore
  \[
    \theta_d
    =
    A_0\mu_d+B_0\cov(\tau_d,X)
    \propto
    \frac{A_0}{B_0}\mu_d+
    \cov(\tau_d,X),
  \]
  with \(\overline W_1\defeq A_0/B_0>0\).
  Substituting back \(X=\ATE_3\) proves the claimed special case.
\end{proof}

\begin{corollary}[Named small-budget coefficient decompositions under an elliptical prior]
  \label{cor:named-objective-elliptical-decompositions-appendix}
  Assume the prior distribution of \(\bm\tau\) is elliptical with finite second moments.
  Then, for every type \(d\), the four named small-budget coefficients satisfy
  \begin{align*}
    \theta_d^{\singlesuccess}
    &=
    \mu_d, \\
    \theta_d^{\twosuccess}
    &=
    \P(\Test_2=1)\,\mu_d
    +
    W_2^{\twosuccess}\,\cov(\tau_d,\ATE_2), \\
    \theta_d^{\singleimpact}
    &=
    \E[\ATE_3] \, \mu_d+
    \cov(\tau_d,\ATE_3), \\
    \theta_d^{\twoimpact}
    &=
    W_1\,\mu_d
    +
    W_2\,\cov(\tau_d,\ATE_2)
    +
    W_3\,\cov(\tau_d,\ATE_3)
  \end{align*}
  for some scalars \(W_1,W_2,W_3\in\R\), not depending on \(d\).
  The coefficient \(W_2^{\twosuccess}\) is defined in \Cref{prop:two-stage-success-elliptical-appendix}; in particular, \(W_2^{\twosuccess}>0\) whenever \(\var(\ATE_2)>0\).
\end{corollary}

\begin{proof}
  The single-stage success and single-stage impact formulas are the general-prior identities in \Cref{prop:first-stage-coefficient-formulas-appendix}.
  The two-stage success formula is \Cref{prop:two-stage-success-elliptical-appendix}.
  The two-stage impact formula is \linkedCref{prop:marginal-impact-interpretation-appendix}.
\end{proof}
